\scanpage{47}
% archive sheet number: 46
\mspage{2}
\begingroup\csname tagsleft@true\endcsname
\[
\begin{aligned}
\sigma_3(z)&=-\frac1z & \Delta_3&=\mathbb V(z^2+1)\ \bigl(=\{i,-i\}\\
&&&\text{si }\exists\ \text{racine primitive}\\
&&&\text{de }\unsure{\varepsilon|4}{-1}\ \text{sur }S\bigr)
\end{aligned} \tag*{(5)}
\]
On trouve donc

\uline{Proposition} Loc pour la top. étale, une représentation monomorphique de \(\mathbb D_4=\mathbb F_2\times\mathbb F_2\) dans \(\operatorname{Aut}(\mathbb P^1_S)\) (sur une base \(S\) à car. résid. \(\ne2\)) est donnée par (4) (\deleted{où} \(\sigma_1,\sigma_2\) sont les autos correspondant aux éléments de la base canonique \((e_1,e_2)\) de \(\mathbb D_4\) sur \(\mathbb F_2\)) -- qui donne bien une représentation monomorphique.

Posons
\[
\begin{aligned}
\Delta&=\Delta_1\amalg\Delta_2\amalg\Delta_3\qquad\bigl(=\{0,\infty,1,-1,i,-i\}\\
&\hspace{14em}\text{dans le cas normalisé}\bigr)
\end{aligned} \tag*{(6)}
\]
on appelle \(\Delta\) la \uline{configuration octaédrale} associée à l'opération de \(V\) sur \deleted{\unsure{\(e\)}} \(X\) -- par descente plate, elle est définie comme diviseur étale de degré \(6\) de \(X/S\)) dès qu'on a une opération monomorphique \inserted{(sur chaque fibre)} d'un groupe \(V\) \inserted{\(\simeq\mathbb D_4\)} sur \(X\), ou plus généralement d'un schéma en groupes \(\underline V\) sur \(S\) loc. isom. au schéma en groupes constant. \deleted{\unsure{Alors}} De façon générale, on appelle \uline{configuration octaédrale} dans \(X/S\) un diviseur étale de degré \(6\) de \(X/S\), qui soit
\endgroup
